\renewcommand{\headline}{Penta$\cdot$ゲーム - 日本語}
\renewcommand{\tocent}{日本語}
\renewcommand{\tocline}{日本語のルール}
\renewcommand{\translator}{Claudeによる下訳、要確認}
\renewcommand{\general}{
1分で説明できるのに、何年経っても夢中になれるゲームです。

2人なら20分ほどで終わります。3人か4人なら90分ほどかかることもあります。

サイコロは使いません。5歳以上ならどなたでも楽しめます。

箱の中身:色のついた駒が4$\times$5個、黒いブロックが5個、グレーのブロックが5個、そしてボード。
}

\renewcommand{\choosext}{駒を選ぶ}
\renewcommand{\choosex}{
プレイヤーはそれぞれ5個の駒を持ちます。箱の中には、プレイヤー1人分の駒セットが入っています。

青、赤、白、緑、黄色の駒が1個ずつあります。

駒は、それぞれの色に対応するボードの隅からスタートし、中央にある大きな交点を目指します。到達すると、その駒は盤から取り除かれます。
}

\renewcommand{\setupt}{準備}
\renewcommand{\setup}{
それぞれの駒を、色が対応する外側の輪の大きな隅に置きます。白い駒は白い隅に、青い駒は青い隅に、というように置いていきます。

黒いブロックは、盤の中央にある5つの交点に置きます。

グレーのブロックは、後で使うので中央に置いておきます。
}

\renewcommand{\objectivet}{目的}
\renewcommand{\objective}{
白い駒は中央の白い交点を、青い駒は青い交点を目指します。目的地は常に、出発点の反対側にある、五角形の中心にある大きな色付きの交点です。

先に\emph{3個}の駒を目的地へ動かした人が勝ちです。
}

\renewcommand{\rulest}{ルール}
\renewcommand{\rules}{
自分の駒を1つ選び、星の線か外側の輪に沿って、好きな方向に好きなだけ動かします。

途中の空いている隅や交点では、止まらずに曲がることができます。

ブロックの上も、他の駒の上も、決して飛び越えることはできません。

\myskip

ただし、すでに駒やブロックがある場所へ進むことはできます:

\myskip

そうして黒いブロックのある場所に進んだ場合は、そのブロックを好きな空いている場所に置き直します。

そうして他の駒がある場所に進んだ場合は、その駒と位置を入れ替えます。

\myskip

自分の駒どうしで位置を入れ替えることもできます。

複数の駒がある場所に進んだ場合は、そのうちの1つを選んで入れ替えます。

\myskip

同じ手を2回続けて行うことはできません。

\myskip

目的地に到達した駒は盤から取り除き、中央に置きます。それから、グレーのブロックを1つ取って、好きな空いている場所に置きます。

そうしてグレーのブロックのある場所に進んだ場合は、そのブロックを盤から取り除きます。

グレーのブロックがなくなった場合は、すでに盤上にあるものを動かします。

\myskip

他のプレイヤーによって自分の駒が目的地まで動かされた場合、自分の番が来たらすぐにその駒を取り除かなければなりません。

\myskip

最後のラウンドは最後までプレイします。
}